\documentclass[
	superscriptaddress,
	twocolumn,
	aps, pre
]{revtex4-1}

\usepackage{amsmath}
\usepackage{amssymb}
\usepackage{hyperref}
\usepackage{graphicx}

% This file defines common macros used in many LaTeX files.
% They are collected here for convenience use.

% packages
\usepackage{amsmath}
\usepackage{xcolor}
\usepackage{dsfont} % double stroke font
\usepackage{xspace}
\usepackage{xparse}


% abbrevations - use \mbox instead of \nolinebreak since this works in captions, too
\newcommand{\ie}{\mbox{i.\hspace{0.125em}e.}\@\xspace}
\newcommand{\eg}{\mbox{e.\hspace{0.125em}g.}\@\xspace}
\newcommand{\etal}{\mbox{et al.}\@\xspace}
\newcommand{\kb}{k_\textnormal{B}}
\newcommand{\kB}{k_\textnormal{B}}
\newcommand{\kbT}{k_\textnormal{B}T}
\newcommand{\kBT}{k_\textnormal{B}T}
\newcommand{\const}{\text{const.}}

% colors, which are often used
\definecolor{plot1}{RGB}{6,115,183}
\definecolor{plot2}{RGB}{255,118,0}
\definecolor{plot3}{RGB}{0,169,25}
\definecolor{plot4}{RGB}{230,0,28}
\definecolor{plot5}{RGB}{0,0,0}
\definecolor{plotBlue}{RGB}{6,115,183}
\definecolor{plotOrange}{RGB}{255,118,0}
\definecolor{plotGreen}{RGB}{0,169,25}
\definecolor{plotRed}{RGB}{230,0,28}

% draft commands
\newcommand{\fixme}[1]{\textcolor{red}{FIXME: #1}}
\newcommand{\todo}[1]{\textcolor{red}{TODO: #1}}
\newcommand{\edit}[1]{\textcolor{blue}{#1}}

% references - use \mbox instead of \nolinebreak since this works in captions, too
\newcommand{\Eqref}[1]{\mbox{Eq.\hspace{0.25em}\eqref{#1}}}
\newcommand{\Eqsref}[1]{\mbox{Eqs.\hspace{0.25em}\eqref{#1}}}
\newcommand{\figref}[1]{\mbox{Fig.\hspace{0.25em}\ref{#1}}}
\newcommand{\figsref}[1]{\mbox{Figs.\hspace{0.25em}\ref{#1}}}
\newcommand{\tabref}[1]{\mbox{Tab.\hspace{0.25em}\ref{#1}}}
\newcommand{\refcite}[1]{\mbox{Ref.\hspace{0.25em}\cite{#1}}}

% math constructs
\newcommand{\im}{{\mathbf{i}}}
\newcommand{\diff}{\text{d}}
\newcommand{\pfrac}[2]{\frac{\partial #1}{\partial #2}}
\newcommand{\deltafrac}[2]{\frac{\delta #1}{\delta #2}}
\newcommand{\difffrac}[2]{\frac{\diff #1}{\diff #2}}
\newcommand{\Difffrac}[2]{\frac{\text{D} #1}{\text{D} #2}}
\newcommand{\takenat}[2]{\left.{#1}\right|_{#2}}
\newcommand{\abs}[1]{|{#1}|}
\newcommand{\Abs}[1]{\left| {#1} \right|}
\newcommand{\norm}[1]{\|{#1}\|}
\newcommand{\Norm}[1]{\left\| {#1} \right\|}
\newcommand{\order}[1]{\mathcal{O}(#1)}
\newcommand{\Order}[1]{\mathcal{O}\left(#1\right)}
\newcommand{\set}[1]{\{{#1}\}}
\newcommand{\Set}[1]{\left\{{#1}\right\}}
\newcommand{\unity}{\mathds{1}}
\newcommand{\identity}{\mathds{1}}
\newcommand{\transpose}[1]{{#1}^\top}
\newcommand{\scalarproduct}[2]{\langle #1 ,\, #2 \rangle}
\newcommand{\Scalarproduct}[2]{\left\langle #1 ,\, #2 \right\rangle}

% math font sizes
\newcommand{\msmall}[1]{\mbox{\small $#1$}}
\newcommand{\mscriptsize}[1]{\mbox{\scriptsize $#1$}}
\newcommand{\mtiny}[1]{\mbox{\tiny $#1$}}

%\newcommand{smallequation}[1]{\begingroup\small#1\endgroup}
\newenvironment{smallequation}{\begingroup\small\ignorespaces}{\endgroup\ignorespacesafterend}

% math sets
\newcommand{\natnum}{\mathds{N}}
\newcommand{\integers}{\mathds{Z}}
\newcommand{\rationals}{\mathds{Q}}
\newcommand{\reals}{\mathds{R}}
\newcommand{\complexes}{\mathds{C}}

% math operators
\DeclareMathOperator{\arctanh}{arctanh}
\DeclareMathOperator{\sech}{sech}
\DeclareMathOperator{\sgn}{sgn}
%\DeclareMathOperator{\max}{max}
%\DeclareMathOperator{\min}{min}
\DeclareMathOperator{\std}{std}
\DeclareMathOperator{\var}{var}
\DeclareMathOperator{\cov}{cov}
\DeclareMathOperator{\CV}{Cv}
\DeclareMathOperator{\STD}{STD}
\DeclareMathOperator{\SEM}{SEM}
\DeclareMathOperator{\erf}{erf}
\DeclareMathOperator{\erfc}{erfc}
\newcommand{\mean}[1]{\langle #1 \rangle}
\newcommand{\Mean}[1]{\left\langle #1 \right\rangle}
\renewcommand\Re{\operatorname{\mathfrak Re}}
\renewcommand\Im{\operatorname{\mathfrak Im}}
\newcommand{\Floor}[1]{\left\lfloor #1 \right\rfloor}
\newcommand{\floor}[1]{\lfloor #1 \rfloor}
\newcommand{\Ceil}[1]{\left\lceil #1 \right\rceil}
\newcommand{\ceil}[1]{\lceil #1 \rceil}
\newcommand{\Nabla}{\vect\nabla}

% special words
\newcommand{\gene}[1]{\mbox{\textit{#1}}}

% special symbols
\DeclareFontFamily{U}{mathx}{\hyphenchar\font45}
\DeclareFontShape{U}{mathx}{m}{n}{<-> mathx10}{}
\DeclareSymbolFont{mathx}{U}{mathx}{m}{n}
\DeclareMathAccent{\widebar}{0}{mathx}{"73}
\newcommand{\vect}{\boldsymbol}
\newcommand{\mat}[1]{\underline{\underline{#1}}}

% symbols for perturbation analysis
\newcommand{\eps}{\varepsilon}
\newcommand{\pert}{\widehat}
\newcommand{\steady}{\widebar}
\newcommand{\pprime}{\prime\prime}
\newcommand{\ppprime}{\prime\prime\prime}

% roman numbers with bars
\def\barroman#1{\sbox0{#1}\dimen0=\dimexpr\wd0+1pt\relax
  \makebox[\dimen0]{\rlap{\vrule width\dimen0 height 0.06ex depth 0.06ex}%
    \rlap{\vrule width\dimen0 height\dimexpr\ht0+0.03ex\relax 
            depth\dimexpr-\ht0+0.09ex\relax}%
    \kern.5pt#1\kern.5pt}}
 
% new environments
\ExplSyntaxOn
\NewDocumentEnvironment{salign} {o}
{\subequations
      \IfNoValueTF{#1} {}{\label{#1}}
      \align}
{\endalign\endsubequations}
\ExplSyntaxOff

% hyphenation

\hyphenation{non-equilibrium}



\renewcommand{\vr}{\vect r}
\newcommand{\vq}{\vect q}
\newcommand{\dc}{\delta c}
\newcommand{\dtc}{\delta \tilde c}

\begin{document}

\title{Linear stability analysis in `py-pde`}
\author{David Zwicker}
\date{\today}

\begin{abstract}
\end{abstract}


\maketitle
\tableofcontents


\section{General derivation}
We consider dynamical equations of the form
\begin{align}
	\partial_t c = f(c, \nabla c, \nabla^2 c, \ldots)
	\;,
\end{align}
which we want to analyze around a homogeneous stationary state, $c(\vr, t) = c_0 + \eps\dc(\vr, t)$.
Here, $c_0$ is chosen such that it is a stationary state, $f(c_0, 0, \ldots)=0$.
To do this analysis, we need to linearize and move to Fourier space.
Below, we check whether the order of operations matters.

\subsection{Linearization first}
Developing the equation up to first order in $\eps$, we thus find
\begin{align}
	\partial_t \dc = \sum_{n=0}^\infty \takenat{\pfrac{f}{(\nabla^n c)}}{c_0} \cdot \nabla^n \dc
	\;.
\end{align}
Note that the term containing derivatives of $f$ is a scalar, vector, or tensor independent of position $\vr$, which is contracted with (derivatives of) the field $\dc$.
Because $\dc$ enters linearly, we impose $\dc(\vr, t) = \exp(\omega t + i \vq.\vr)$, effectively moving to Fourier space,
\begin{align}
	\label{eqn:growth_rate_goal}
	\omega(\vq) = \sum_{n=0}^\infty \takenat{\pfrac{f}{(\nabla^n c)}}{c_0}\cdot (i\vect q)^n
	\;,
\end{align}
which should be the equation we want.

\subsection{Fourier transform first}
Since the form discussed above requires derivatives of the right hand side with respect to the derivatives of fields, which is difficult to implement in \texttt{sympy}, we here check whether we can first replace the derivatives by going to Fourier space and then linearize.
We thus first consider the Fourier transformed function using $c(\vr, t) = \tilde c(\vq, t) e^{i\vr\vq}$,
\begin{align}
	\partial_t \tilde c = \tilde f(\tilde c, i\vq \tilde c, -\vq^2 \tilde c, \ldots) = \tilde g(\tilde c, \vect q)
	\;.
\end{align}
We then expand the field up to first order in $\eps$ using $\tilde c(\vq, t) = c_0 + \eps \dtc(\vq)e^{\omega t}$,
\begin{align}
	\label{eqn:growth_rate_hypothesis}
	\omega(\vq) &= \takenat{\pfrac{\tilde g}{\tilde c}}{c_0} 
\notag\\
	&= \sum_{n=0}^\infty \takenat{\pfrac{\tilde f}{((i\vq)^n\tilde c)}}{c_0} \cdot (i\vq)^n
	\;,
\end{align}
which is identical to \Eqref{eqn:growth_rate_goal}.

\subsection{Rule-of-thumb}
What about simply replacing every $\nabla$ by $i\vq$?
Does this work for functions where the differential operators are not directly applied to $c$?


\section{Concrete examples}

\subsubsection{Compound functions}
To see whether \Eqref{eqn:growth_rate_goal} and \Eqref{eqn:growth_rate_hypothesis} agree, we test a simple example,
\begin{align}
	f &= \nabla\bigl[A(c) \nabla B(c)\bigr]
\notag\\
	 &=(\nabla A(c)). (\nabla B(c))  + A(c) \nabla^2 B(c)
\notag\\
	 &= A'(c)B'(c) (\nabla c)^2  + A(c)B''(c) (\nabla c)^2 + A(c)B'(c) \nabla^2 c
	\;.
\end{align}
First, \Eqref{eqn:growth_rate_goal} can be brought to the form where the terms are derivatives are directly applied to $c$.
\begin{align}
	\omega(\vq) &= \sum_{n=0}^\infty \takenat{\pfrac{f}{(\nabla^n c)}}{c_0}\cdot (i\vect q)^n
%\notag\\
%	&= \bigl[A''(c_0)B'(c_0) + A'(c_0)B''(c_0)\bigr](\nabla c_0)^2  + \bigl[A'(c_0)B''(c_0) + A(c_0)B'''(c_0)\bigr] \nabla^2 c_0+
%	2 i A'(c_0)B'(c_0)\vq \nabla c_0 - A(c_0)B''(c_0) \vq^2  % plus extra terms
\notag\\
	&= - A(c_0)B'(c_0) \vq^2
	\;.
\end{align}
Using the rule-of-thumb, we obtain $\omega(\vq) = i\vq \bigl[A(c) i\vq B(c)\bigr]$, which gives the right scaling with $\vq$, but does not lead to any derivatives in the functions $A$ and $B$.
\todo{Apparently, we need to linearize in the field first and then perform the Fourier transform. We thus need to tell \texttt{sympy} how to linearize the differential operators}
This would imply
\begin{align}
	f &= \nabla\bigl[A(c_0 + \dc) \nabla B(c_0 + \dc)\bigr]
\notag\\ & \approx
	\nabla\bigl[\bigl(A(c_0) + A'(c_0)\dc\bigr) \nabla\bigl(B(c_0) + B'(c_0) \dc\bigr)\bigr]
\notag\\ & \overset{\text{nabla const}}{\approx}
	\nabla\bigl[\bigl(A(c_0) + A'(c_0)\dc\bigr) B'(c_0)  \nabla \dc\bigr]
\notag\\ & \overset{\text{linearize}}{\approx}
	\nabla\bigl[A(c_0) \nabla B'(c_0) \dc\bigr]
\notag\\ & \overset{\text{collect}}{\approx}
	A(c_0) B'(c_0) \nabla^2 \dc
\notag\\ & \overset{\text{Fourier}}{\approx}
	A(c_0) B'(c_0) (-\vq^2 \dc)
\end{align}
This works well in this simple case, but we somehow need to tell sympy how to simplify the expressions.

\subsection{Compositions}
For the simple case
\begin{align}
	f = A(\nabla^2 c)
\end{align}
we find
\begin{align}
	f = A(c_0 + \nabla^2 \delta c)
	\;.
\end{align}
Assuming $\nabla^2 \delta c$ is also small, we have
\begin{align}
	f \approx A(c_0) + A'(c_0) \nabla^2 \delta c
\end{align}
For the more complicated case
\begin{align}
	f = B\bigl(\nabla^2\bigl[A(\nabla^2 c)\bigr]\bigr)
\end{align}		
we find
\begin{align}
	f &= B\bigl(\nabla^2\bigl[A(c_0 + \nabla^2 \delta c)\bigr]\bigr)
\notag\\
	& \approx B\bigl(\nabla^2\bigl[A(c_0) + A'(c_0) \nabla^2 \delta c\bigr]\bigr)
\notag\\
	& = B\bigl(A'(c_0) \nabla^4 \delta c\bigr)
\notag\\
	& \approx B'(0) A'(c_0) \nabla^4 \delta c
\end{align}
\todo{Can we simply replace all fields by $c_0 + \epsilon\delta c$ and use sympy to expand in $\epsilon$? I guess we need to tell sympy that the differential operators are linear, but this might be possible. We might need to do this from the inside out, ensuring that each leaf-expression is linear. We can do this for each linear operator (expand inside). During this linearization we might be able to replace the differential operators by their Fourier transform }


%\bibliographystyle{apsrev4-1}
%\bibliography{bibdesk.bib}

\end{document}  